% !TEX root = ../fade.tex

\subsection{Motivation}

	There are many decades of literature regarding the problem of estimating a mapping $\mathcal{P}\to\mathcal{O}$, from some parameter space $\mathcal{P}$ to an observable space $\mathcal{O}$, conditioned on a set of observed data $\vec{T} = \{(\vec{y} \in \mathcal{O}, \vec{x}\in \mathcal{P})\}$. The vast majority of this work is on `point predictors', wherein one is able to reliably compute $\mathbb{E}(\hat{\vec{y}} | \vec{w}, \vec{T})$ for some new proposed set of covatiates $\vec{w}\in \mathcal{P} \backslash{} \vec{T}$.

	More modern methods, such as Gaussian Process Emulation, are able to also predict a distribution, but in doing so must assume a fixed parametric family which is attached to the point predictor.

	For our purposes however, we are more interested in a \textit{distributional predictor}, which would enable us to predict not just the behaviour of the mean, but the entire distribution, $p(\ny| \nx, \vec{T})$.

	Whilst there are many methods which allow us to \textit{quantify the uncertainty} of a prediction, there are few which allow us to fully reconstruct an accurate \textit{distribution}, if the underlying `truth' that is being distributed is genuinely a stochastic process, rather than a `deterministic process with unknown noise'.


\subsection{The Data}

	We supposes that there exists some observable $\y \in \R^q$ which may be extracted from a function $\hat{M}$ evaluated at a parameter vector $\vec{P} \in \vec{\Pi}$:
	\begin{equation}
		\y = \hat{M}(\vec{P})
	\end{equation}


	For the cases that interest us, $\hat{M}$ is a simulational (though potentially stochastic) model of a system where the input dimension, $\text{dim}(\vec{\Pi})$ is large.

	\subsubsection{Parameter Promotion}

		For the purposes of a statistical emulation, we suppose that this full parameter space may be divided into two subspaces:

		\begin{enumerate}
			\item \textbf{Conditional Variables} These are the covariates of interest for the final emulator.
			      They define the \emph{emulation space}
			      \begin{equation}
				      \mathcal{E} = \mathbb{R}^n
			      \end{equation}
			      and are denoted by $\vec{x}$.

			\item \textbf{Latent Variables} The remaining parameters of the mechanistic model.
			      These variables are not explicitly represented in the emulator and instead contribute stochastic variability to the output distribution. They are denoted $\vec{z}$, and we assume that there exists some prior on their values, which may be conditioned on the covariates $P(\vec{z} | \vec{x})$.
		\end{enumerate}



	\subsubsection*{Why use manual promotion?}

		PCA and related nonlinear projection methods provide a method in which the machine computes a projection into the space of the `most important variables':
		\begin{equation}
			\hat{\vec{\chi}} = \mathcal{T}(\vec{P})
		\end{equation}
		It would then be possible to learn $p(\y | \hat{\vec{\chi}})$, which is an easier task since $\hat{\chi}$ is a lower-dimensional entity, and the projection has (in theory) captured the bulk of the underlying correlations.

		However, this approach requires specifying a full parameter vector $\vec{P}$ whenever the model is used to generate a prediction. For complex mechanistic systems this is often undesirable. For example, in the mosquito model we wish to predict outcomes conditioned on the gene drive parameters, without needing to specify rainfall realisations and mosquito lifetimes.

		In short, we wish to simulataneously marginalise over these parameters, and fit the covariates of interest.

\subsection{Generating Data}


	\subsubsection{Latent Variables \& Priors}

		The `obvious' solution for generating the data we need for our emulator would be to sample $\vec{x}$ uniformly from $\mathcal{E}$, and $\vec{z}$ from the latent prior, $P(\vec{z})$, and then compute $\vec{y} = \hat{M}(\x, \vec{z})$. Our dataset would then be $\D =  \{(\vec{x}_i, \y _i)\}_{i=1}^{D}$.

		However, this would have the major drawback that if -- at a later date -- our priors on the latent variables were updated, then we would need to completely regenerate the training data. This would be a significant bottleneck in computation time: for the models of interest to us, the vast majority of the computational expense will be in creating $\D$.

		Instead, we suppose that we maintain two datasets, where the latent priors are drawn from $P_\text{sample}(\vec{z},\vec{x})$, which is not necessarily the same as our current prior on $\vec{z}$.
		\begin{equation}
			\D = \{(\vec{x}_i, \y_i, \eta_i)\}_{i=1}^{D}~~~\quad~~~\mathcal{Z} = \{ \vec{z}_i, \zeta_i \}_{i=1}^D
		\end{equation}
		The ordering between these sets is the same (so $\y_i = \hat{M}(\x_i, \vec{z}_i)$ is true). The parameters $\eta_i$ and $\zeta_i$ are defined as:
		\begin{spalign}
			\eta_i & = P_\text{sample}(\vec{z}_i)
			\\
			\zeta_i & = \frac{P_\text{latent}(\vec{z}_i)}{\eta_i} \label{E:DefineWeights}
		\end{spalign}
		Where  $P_\text{latent}$ is the \textit{current prior on the latent variables}, and we recall $P_\text{sample}$ is the distribution used to generate $\{\vec{z}\}$. This means that $\D$ is immutable (unless more data is added), and changes to the prior require simply recomputing $\zeta_i$ - a trivial task.

		When computing the log-posterior conditioned on the observed evidence, we simply need to correct by a factor $\zeta_i$:

		\begin{equation}
			\mathcal{L}(\D, \mathcal{Z}) = \sum_{i=1}^D \zeta_i \times \ln p(\y_i | \x_i,....)
		\end{equation}


		We can see that if $\zeta_i = 1$ (which occurs when $P_\text{sample} = P_\text{latent}$), then this reduces to the obvious case. Otherwise, it means that we have an effective sample size
		\begin{equation}
			D^* = \frac{\left(\sum_i \zeta_i \right)^2}{\sum_i \zeta_i^2}
		\end{equation}

		The choice of sampling distribution must therefore be informed on a best-guess understanding of how the latent prior may evolve to ensure that $D^*$ is a reasonable size for both the current latent prior, and any future updates. If $D^* \ll D$, then it will be necessary to generate more data, which is to be avoided under the current paradigm.
	\subsubsection{x-Clustering}

		We also note that clustering the data at values of identical $\vec{x}$ provides both a vast computational speed up, and negates a danger posed by the mean-fitting discussed in \S\ref{S:BLP}. That is, instead of generating $D$ unique values of $\vec{x}$ (and then $\vec{z}$ and hence $\y$), we suppose that only $T$ values of $\x$ are unique.

		At each unique $\x_t$, we then generate $A_t$ pairs of ($\vec{z}$, $\vec{y}$) pairs\footnote{Indeed, if the model is stochastic, it may be that we keep $\vec{z}$ constant and sample several $\y$ from a single $(\x_t, \vec{z}_a)$ pair - the optimal distribution here is to be determined by the trainer.}, to form a clustered dataset

		\begin{equation}
			\mathcal{D} = (\left\{ \x_t, \{\y_t^a, \eta_t^a\}_{a=1}^{A_t} \right\}_{t = 1}^{T}
		\end{equation}

		This provides a speed up since the bulk of the computation time during training is spent computing $W_i(\x)$ - if this value can be reused $\approx A_t$ times, it provides a commensurate $\approx A_t\times$-speedup.

